\documentclass[11pt]{article}
\usepackage[margin=1.1in]{geometry}
\usepackage{amsmath,amssymb,url}
\newcommand{\grad}{\nabla}
\newcommand{\eps}{\varepsilon}
\setlength{\parskip}{4pt}

\title{\vspace{-2em}Open Letter to the Clay Mathematics Institute\\[0.3em]
\large On the Formulation and Solution of the Navier--Stokes Problem}
\author{Claes Johnson\\[0.2em]
\small KTH Royal Institute of Technology, SE-100 44 Stockholm, Sweden\\
\small \texttt{claesjohnson@gmail.com}}
\date{\today\\[0.4em]\small\textit{written up with AI assistance (Claude, Anthropic)}}

\begin{document}
\maketitle

\noindent
To the Board of the Clay Mathematics Institute,

\medskip

I write about the Navier--Stokes Millennium Problem: to propose a reformulation of it, to present a
solution of the problem so reformulated, and to explain why I believe the reformulation is necessary
rather than evasive. None of this is new. It was published in 2007, in a chapter of a book written
with Johan Hoffman, and it has been put to the Institute before. What is new is the occasion. On
8~September 2026 OpenAI announced a solution of the problem as officially posed~\cite{openai}, and
that announcement makes it worth setting out, once more and in one place, what I take the meaningful
question to be, what answering it consists of, and what the announced solution does and does not
settle.

I should say at the outset where I am writing from, because the proposal did not arrive from outside
the subject. It is the end of a long piece of work, and the reformulation was reached by computing,
not by reflecting on the problem statement.


\section{Where this comes from: a computational programme, 1970s--2026}

\paragraph{The 1970s: finite elements for elliptic problems.} My work began with the finite element
method for elliptic partial differential equations --- error analysis, a priori and a posteriori
estimates, and the question that has organised everything since: not whether a method converges in
principle, but how large the error is in \emph{this} computation, on \emph{this} mesh, in a quantity
somebody actually wants. For elliptic problems that question has a clean answer, and the answer is
obtained by duality: the error in a functional of the solution is the residual of the computed
solution weighted by the solution of an adjoint problem. Everything below is that idea, carried as
far as it will go.

\paragraph{The 1980s: hyperbolic problems, and fluid flow.} The methods were then extended to
hyperbolic and convection-dominated problems, where the difficulties are of a different order:
solutions develop sharp layers and discontinuities, Galerkin's method loses stability, and
stabilisation has to be supplied without destroying the error control. Streamline diffusion and
least-squares stabilisation of the residual came out of that period, and with them the first
computations of fluid flow in which the dissipation was generated by the scheme in proportion to the
residual rather than prescribed by the person running it.

\paragraph{The 1990s and 2000s: turbulent flow, adaptivity, error control.} Carried to the
Navier--Stokes equations at high Reynolds number, that programme produced the result which forced
the reformulation. Computed turbulent solutions do not converge pointwise under mesh refinement ---
two meshes give two velocity fields that differ everywhere and go on differing --- while
\emph{mean-value outputs} such as drag and lift converge to a few per cent, with the error bounded
by duality against a computed adjoint solution rather than estimated by comparing successive meshes.
That is not a defect of the computations. It is what the equations do, and once one has seen it
repeatedly it is no longer possible to regard pointwise regularity as the property a formulation
should ask about.

\paragraph{2007: the book, and the chapter.} This was set out in \emph{Computational Turbulent
Incompressible Flow}~\cite{cfd}, written with Johan Hoffman and published by Springer in 2007 as
volume~4 of \emph{Applied Mathematics: Body and Soul}. One chapter --- Chapter~14 --- is given to the
Clay problem: it argues that the official formulation asks about the wrong property, proposes the
reformulation in terms of output wellposedness, and presents computed solutions as its answer. What
follows in this letter is that chapter's content, recalled and brought up to date. The earlier
statement of the approach is in~\cite{newapproach}, the residual-based dissipation in~\cite{hjv}, the
parameter-free adaptive formulation in~\cite{hjja}, the resolution of d'Alembert's paradox
in~\cite{dalembert}, and a theory of flight built on the same computations in~\cite{flight}.

\paragraph{2007--2026: compressible flow, and where the programme now stands.} The development since
has gone to compressible flow and to aeronautical Reynolds numbers, and it has been carried
substantially by Johan Jansson. The adaptive methodology was built into software --- Unicorn, for
adaptive simulation of turbulent flow and fluid--structure interaction~\cite{unicorn}, within a
framework for massively parallel adaptive finite element computational fluid
dynamics~\cite{massparallel}, and then FEniCS-HPC, for automated predictive high-performance finite
element computing with aerodynamics applications~\cite{fenicshpc} and coupled
multiphysics~\cite{fenicsmulti}, with parallel-in-time computation of turbulent flow
in~\cite{paralleltime} --- and applied at aeronautical Reynolds numbers in the DLR-F11 aircraft
case, time-resolved and adaptively refined~\cite{dlrf11}, and in the parameter-free formulation
of~\cite{hjja}. The methodology now takes a leading role in the fluid mechanics workshops where
computational predictions are compared against measurement. The argument of this letter is therefore not a position held in isolation. It is the
reading of the Navier--Stokes equations that a working computational programme has arrived at, and
the programme has been tested in the ordinary way, against experiment, for two decades.


\section{The reformulation}

\begin{quote}
\textbf{Problem.} \emph{Constructive existence and qualitative uniqueness of weak solutions to the
incompressible Navier--Stokes equations with small viscosity, that is at high Reynolds number.}
\emph{Given data:} a domain $\Omega\subset\mathbb{R}^3$, initial data $u_0$ on $\Omega$, a forcing
$f$ on $\Omega\times(0,T)$, and boundary conditions on $\partial\Omega$.
\end{quote}

Both adjectives carry weight, and each is weaker than the classical demand in one respect and
stronger in another. In both the trade is the same: less is asked of the pointwise field, more of
what can be done with it.

\paragraph{Constructive existence, as against existence only.} Classical existence asserts that the
solution set is non-empty. Leray's theorem~\cite{leray,lions} is of that kind: obtained by
compactness, it establishes that a weak solution exists and says nothing about which one, what it
looks like, or how to reach it. For \emph{small} data --- the low Reynolds number case for suitably
scaled problems --- global existence of a \emph{smooth} solution has been known for sixty
years~\cite{fujitakato,kato}, so the difficulty lies entirely at the other end, which is the end the
reformulation addresses. For a problem whose whole difficulty is that the solutions are complicated, that is
close to no information. I ask instead that existence be established \emph{by exhibition}, and
concretely that five things be produced rather than asserted: the mesh, with its cell count and
smallest cell; the discretisation carrying the solution, here continuous piecewise linear in space
and time with least-squares control of the residual; the residual itself, measured in the norm the
error bound uses; the computed stability factor for each output; and then, for each output --- the
drag of a body, the lift of a wing, a dissipation rate, a pressure distribution --- a value with a
bound on its error. That is what the solution de facto \emph{is}. Ask a non-constructive existence
theorem for the drag and there is nothing there to ask.

There is more to this than error control. A computed solution is an elaborate explicit construction
of a velocity--pressure field, some millions of numbers on a spacetime mesh, every one available ---
and such an object can be taken apart. One can ask it where vorticity concentrates and in what
shape, how the velocity gradient decomposes at a point, where energy passes from large scales to
small and through what structures, where the dissipation actually sits as against where a model would
put it. None of these questions can be put to a solution known only to exist. A solution concept
whose objects can be inspected generates conjectures; one whose objects cannot be exhibited does not.
The triple decomposition of the velocity gradient~\cite{triple}, and the energy stability analysis
built on it~\cite{energystab}, came out of looking closely at resolved velocity fields, and at
measured ones~\cite{musci}.

\paragraph{Qualitative uniqueness, as against uniqueness only.} Classical uniqueness asks that the
data determine the field pointwise. At high Reynolds number this is false, and not marginally. What
the data \emph{do} determine is the qualitative content, and concretely this is what two computations
on independently generated meshes are found to share: the drag agrees to a few per
cent~\cite{drag} and the lift likewise; the separation line sits in the same place; the wake is
organised the same way and sheds at the same frequency; the total dissipation rate is the same. What
they do not share is the velocity, which at any chosen point differs by order one and goes on
differing under refinement.

Qualitative uniqueness is the demand that this be so, and that the line between the two lists be
drawn \emph{quantitatively} rather than by assertion. That last clause is what keeps it from being a
soft version of uniqueness. The adjoint solution attaches to each functional a computed stability
factor: where it is moderate the output is determined by the data to a stated precision, and
refinement where the adjoint is large will improve it; where the factor is large the output is not
determined, and no refinement will make it so. The claim is falsifiable one functional at a time, and
a functional can fail it.

Three further things distinguish this from the official statement. \emph{The domain is given}, not
$\mathbb{R}^3$ with decay --- real flows occupy bounded domains, and the boundary is where drag and
lift are produced. \emph{Small viscosity is part of the problem}, not a parameter fixed in advance:
$\nu\to0$ is the limit of interest and $T$ is finite, which is the difference between asking about
one Reynolds number at a time and asking about turbulence. And \emph{weak solutions are asked for,
smoothness is not} --- by the Onsager theorem~\cite{onsager,isett} non-smoothness is not a defect to
be removed but the mechanism by which a turbulent flow dissipates energy at a rate independent of
viscosity --- the dissipation anomaly, or zeroth law, observed since Taylor~\cite{taylor} and still
under active examination~\cite{iyer}. Requiring smoothness excludes the solutions that occur.


\section{The solution}

\begin{quote}
\textbf{Solution.} \emph{Computed solutions, as a form of constructed solution, for a representative
choice of specific data --- constructive existence by exhibition, qualitative uniqueness by the
adjoint.} For given $\Omega$, $u_0$, $f$ and boundary conditions, exhibit a computed weak solution
together with quantitative control of the outputs: a residual small in the relevant norm, an error
bound for each output obtained by duality from a computed solution of the adjoint equations, and a
demonstration that the output converges under refinement while the velocity field does not.
\end{quote}

\paragraph{Wellposedness is a property of outputs, not of solutions.} This is the point on which the
reformulation turns. Hadamard's criterion has three parts --- existence, uniqueness, continuous
dependence on the data --- and the third carries the physical content, because data are never exact.
For a turbulent solution the pointwise velocity is not determined: perturb the data in the last
decimal and the field is different everywhere within a short time. In that sense the solution is not
wellposed and never will be. But nobody measures the pointwise velocity of a turbulent flow, and no
engineering question depends on it. What is measured are mean-value outputs,
\begin{equation}\label{eq:output}
M(u) = \int_Q u\cdot\psi\,dx\,dt ,
\end{equation}
and the question of wellposedness must be asked of \emph{them}. The claim is that $M(u)$ is stable
for turbulent solutions even though $u$ is not.

\paragraph{The dual solution is a weight function connecting residual to output.} Output quality is
not assessed by inspection; it is computed. The error in an output is, to leading order, the residual
of the computed solution tested against a dual weight function $\varphi$,
\begin{equation}\label{eq:dualrep}
M(\hat u) - M(\hat U) \;\approx\; \int_0^T\!\!\int_\Omega R(\hat U)\cdot\varphi \,dx\,dt ,
\end{equation}
so $\varphi$ is what converts a residual --- which one has, and can measure --- into an error in the
quantity one cares about. Where $\varphi$ is large the residual matters and the mesh must be refined;
where it is small the residual is harmless however large; and its size overall is the stability
factor that says whether the output is determined at all. This yields an a posteriori bound
\begin{equation}\label{eq:apost}
|M(u) - M(U)| \;\le\; \int_Q R(U,P)\cdot\varphi\,dx\,dt \;\le\; S\,\|R(U,P)\|_{-1},
\end{equation}
in which $S$ is not assumed but computed.

And $\varphi$ is obtained by solving a definite problem: the Navier--Stokes equations linearised
about the computed solution, run \emph{backward in time} from the output functional as data at
$t=T$. This is the same object, and the same computation, as the backward pass in the training of a
neural network --- the adjoint of a forward evolution, propagating sensitivity from an output to
everything that fed it. Fluid mechanics and machine learning are doing the same linear algebra in
opposite directions through the same trajectory.

\paragraph{Why the backward problem is stable for a turbulent solution.} That the dual problem can be
solved at all is the crux, and it is where intuition points the wrong way. A linearisation about a
chaotic trajectory has exponentially growing modes; run backward, one expects $\varphi$ to blow up
and the whole assessment to be vacuous. It does not, and the reason is the oscillation. A turbulent
solution does not grow exponentially; it \emph{oscillates}, and it oscillates precisely \emph{because}
it is the result of exponential growth followed by exponential decay --- a structure is stretched and
amplified, becomes unstable, breaks down and dissipates, and the next takes its place. The linearised
coefficients along such a trajectory change sign and structure repeatedly, so the growth $\varphi$
accumulates over one interval is largely undone over the next: what would be a monotone exponential
over a smooth field becomes a product of factors alternately greater and less than one. Over a
\emph{smooth} field there is nothing to cancel, and it grows.

So turbulence stabilises the outputs that smoothness destabilises, and by the same mechanism that
makes it unpredictable pointwise. Smooth solutions are not wellposed. Turbulent solutions are
wellposed in mean-value outputs. The physically meaningful solution concept is the one the official
formulation treats as the failure case.

\paragraph{Turbulence is a necessity, not a possibility.} It is worth adding why the third case is
not one option among several. High Reynolds number flow starting out laminar develops shear
instabilities with sharp velocity gradients, which are eventually damped by viscosity in turbulent
flow. All high Reynolds number laminar flow is thus unstable and cannot continue to exist over time;
it turns turbulent in order to continue to exist. Turbulence is what the equations must do.


\section{Cases of computed solutions, with assessment of uniqueness}

The method is General Galerkin, G2: a stabilised finite element method of Galerkin least-squares
type~\cite{fehn}, continuous piecewise linear in space and time, $\mathrm{cG}(1)\mathrm{cG}(1)$, with least-squares control of the residual. The
stabilisation is not a subgrid model --- no eddy viscosity is prescribed, no wall function, no
closure, no constant fitted to measurement. The dissipation is whatever the residual makes it, which
means it is generated by the flow being computed rather than supplied by the person computing it.
That matters for the present argument specifically: a computation carrying a fitted turbulence model
could not be used to say what Navier--Stokes contains, because the turbulence would have been put in
by hand.

Across the cases~\cite{cfd,hjv,hjja,newapproach,dlrf11} the pattern is the one the reformulation
predicts. The velocity field fails to converge under refinement and is not expected to. The
mean-value output converges to a few per cent, and the error in it is bounded by~\eqref{eq:apost},
not estimated after the fact. Mean drag and mean lift for turbulent flow past a bluff body come out
to a few per cent on meshes of a few hundred thousand points --- orders of magnitude coarser than a
direct simulation resolving the Kolmogorov scale would require. And the adjoint supplies the
assessment of uniqueness directly: where its stability factors are moderate the output is unique in
the sense that matters; where they are large the output is not predictable and refinement will not
make it so. That last case is not a failure of the method but an answer --- the equations do not
determine that quantity, and no computation should be believed which claims otherwise. Further
material is collected under the label \emph{clay problem}~\cite{blog}.


\section{The official formulation recalled}

The problem as set out by Fefferman~\cite{fefferman} concerns
\begin{equation}\label{eq:ns}
\dot u + (u\cdot\grad)u - \nu\Delta u + \grad p = f,\qquad \grad\cdot u = 0,
\end{equation}
on $\mathbb{R}^3$ or a periodic box, with $u(\cdot,0)=u^0$ smooth and divergence free, and asks
whether solutions remain smooth for all time or break down. I recall it here rather than at the
start, because a formulation is best judged against one that has a solution.

\paragraph{The third case has no name in the formulation.} Turbulent flow is not smooth in any useful
sense: the gradients scale as
\begin{equation}\label{eq:grad}
|\grad u| \sim \nu^{-1/2},
\end{equation}
so at $\nu=10^{-6}$ they are of order $10^{3}$ times the velocity scale. Nor is it singular: the
velocity stays bounded. The two halves are one scaling seen from either end,
\begin{equation}\label{eq:delta}
\eta \sim \Big(\frac{\nu^3}{\eps}\Big)^{1/4} \sim \nu^{3/4},
\qquad
\delta u(\eta) \sim (\eps\,\eta)^{1/3} \sim \nu^{1/4},
\qquad
\frac{\delta u(\eta)}{\eta} \sim \nu^{-1/2},
\end{equation}
the first two being the Kolmogorov--Obukhov scaling~\cite{kolmogorov41,obukhov41} --- refined by the
intermittency corrections of~\cite{kolmogorov62,obukhov62} and the multifractal picture
of~\cite{frischparisi} --- and the middle relation being H\"older continuity with exponent $1/3$, the
Onsager exponent. As
$\nu\to0$ the structures refine without limit while the velocity stays bounded at every scale: the
gradient diverges, so the solution is not smooth, and the velocity does not, so it is not singular.
A H\"older-$1/3$ field is exactly the object that is neither differentiable nor unbounded, which is
why no amount of regularity theory reaches it. By the Onsager theorem~\cite{onsager,isett}, proved after the partial results
of~\cite{eyink,cet,dls,bdlis}, this category has occupants, and a local energy equation can be written in which the non-smoothness is
the direct cause of the dissipation~\cite{duchonrobert}.

The Institute's own description of the problem opens on turbulence: waves follow the boat across the
lake, turbulent air currents follow the aircraft, and the challenge is to unlock the secrets hidden
in the equations. That is accurate about where the secrets are. The formal statement that follows
contains no turbulence; it asks whether the velocity stays finite. Between the description and the
problem the subject has been sorted out --- not by oversight in the description, which names the
phenomenon correctly, but by the formalisation, which admits no predicate that could distinguish a
turbulent solution from a laminar one.

\paragraph{The alternatives are not complementary.} Fefferman's four statements do not form two pairs
of negations. (A) and (B) are posed with $f\equiv0$ and quantify over \emph{all} admissible data:
does every datum give a global smooth solution? (C) and (D) ask instead for the \emph{existence} of
admissible data \emph{together with an admissible smooth forcing} $f$ for which smoothness fails. So
(C) is not $\neg$(A). A blowup driven by an external force leaves (A) untouched: it does not show
that a fluid tears itself into a singularity, only that something can be done to it that does.


\section{The announced solution, and the viscosity test}

On 8~September 2026 OpenAI announced a proof that a singular solution exists~\cite{openai}, obtained
with ten thousand agents and of the order of $3\times10^{11}$ output tokens~\cite{buckmaster}. I take
it to be a serious piece of work and I do not dispute it. If it stands, the Clay problem as posed is
solved. Two observations follow, and the second is the substance of this letter.

\paragraph{What it is a solution of.} Theorem~1.1 asserts that for every $\nu>0$ there exist a force
$f\in C^\infty_c(\mathbb{R}^3\times(0,\infty))$, a compact set $K$, and smooth fields $u,p$ on
$\mathbb{R}^3\times[0,1)$ solving \eqref{eq:ns} with $u(\cdot,0)=0$, supported in $K$, with
$\sup_t\|u(t)\|_{L^2}<\infty$ and $\limsup_{t\uparrow1}\|u(t)\|_{L^\infty}=\infty$. The flow starts
from rest, so everything in it comes from the force --- and the order of construction is the point.
The velocity and pressure are built first, to have the required growth, and the force is defined as
what is left over, by making its momentum residual the restriction of a force in
$C^\infty_c$. The singularity is therefore supplied to the equations rather than generated by them,
and the admissible class is broad enough to let it be supplied. That the class carries the result is
visible from the other side too: replacing the $C^\infty$ force by a spatially real-analytic one
rules the mechanism out~\cite{civ}. The paper states its own scope in just these terms --- it
establishes alternative (C), with the periodic case giving (D). It does not claim (A) or (B), and by
the previous section it could not.

\paragraph{A shift in what a proof is, and why it bears on the reformulation.} Something larger than
the Navier--Stokes problem is visible here. With AI, mathematics is undergoing a monumental shift from
symbolic analysis towards constructive computing --- from proofs establishing that an object exists to
proofs that exhibit one. The announced solution is an example of it. It does not argue abstractly that
a singular solution must exist: it \emph{builds} one, a velocity and pressure field specified in
advance to have the required growth, with the force then defined as the residual of that construction.
Ten thousand agents and of the order of $3\times10^{11}$ tokens produced an explicit object, not an
existence argument.

That is the direction in which this letter asks the formulation to move, and it is why the
reformulation is not an eccentric demand. Constructive existence is becoming the normal form of a
mathematical result rather than a weaker substitute for the real thing. The difference between the
announced construction and the one I offer is therefore not that one is constructive and the other
abstract --- both are constructions, and both are specific. It is what they construct: a force built
to order to produce a chosen singularity, against a flow computed forwards from data a fluid actually
has, with a bound on what may be asked of it. Both exhibit. Only one exhibits a fluid.

\paragraph{The viscosity test.} The test is general: it applies to any construction that exhibits a
singularity at prescribed forcing. Navier--Stokes rescales --- if $v$ solves the equations at
viscosity one with forcing $g$, then
\begin{equation}\label{eq:rescale}
u(x,t) = a\,v(bx,ct),\qquad a = \nu b,\quad c = ab,\quad f = a^2 b\,g,
\end{equation}
solves them at viscosity $\nu$ --- and the map is onto: \emph{one} unit-viscosity blowup generates
the entire family. A theorem asserting blowup ``for every $\nu>0$'' is therefore not a family of
results. It is one result, rescaled, containing exactly what its $\nu=1$ member contains.

Everything then turns on what \eqref{eq:rescale} does to the constants, and since it is the forcing
class that the official formulation fixes, two normalisations bear on it. Hold the solution at size
one, $a=1$: then $b=c=\nu^{-1}$ and the forcing amplitude is $a^2b=\nu^{-1}$, so the force required
\emph{blows up} as the viscosity falls and leaves any fixed admissible class --- including the one
Fefferman specifies, which fixes its constants once and for all. Or hold the forcing in a fixed
class, $a^2b=1$: then $b=\nu^{-2/3}$ and
\begin{equation}\label{eq:shrink}
\text{amplitude } a = \nu^{1/3}\to 0,\qquad
\text{length } b^{-1} = \nu^{2/3}\to 0,\qquad
\text{time } c^{-1} = \nu^{1/3}\to 0 .
\end{equation}
The singularity shrinks to nothing: at $\nu=10^{-6}$ the velocity reaches $10^{-2}$ on a scale of
$10^{-4}$ before blowing up. Neither branch survives the limit. One exits the admissible data, the
other converges to the zero solution, and there is no third option.

Set this beside the flow the problem is advertised as being about. Turbulence at small viscosity has
bounded velocity at every scale and gradients growing like $\nu^{-1/2}$, realised by
\eqref{eq:delta} on structures of size $\nu^{3/4}$ across which the velocity varies by $\nu^{1/4}$
--- amplitude fixed, structure refining. The constructed family has amplitude vanishing like
$\nu^{1/3}$, which is the opposite behaviour, or else forcing diverging like $\nu^{-1}$, which is not
a fluid at all. A singularity that exists at every fixed $\nu$ and disappears as $\nu\to0$ is not a
feature of high-Reynolds-number flow; it is a feature of a formulation which admits forces no
mechanism supplies.


\section{In closing}

The material above has been before the Institute, in one form or another, for some twenty years. It
was published in 2007 in Chapter~14 of~\cite{cfd}, it rests on a computational record documented in
refereed journals from~\cite{newapproach} onward, and the programme it belongs to has continued since
into compressible flow and aeronautical Reynolds numbers. I have not changed the argument, and
nothing in the intervening twenty years has required me to.

What has changed is the actuality, and it has changed in two ways. The first is the shift described
above: as mathematics moves from symbolic analysis to constructive computing, a formulation asking for
constructive existence and qualitative uniqueness is no longer asking for something unusual. The
second is that a solution to the official formulation has now been announced, and it is of exactly
the kind the formulation invites: a force constructed to produce a chosen
singularity, admissible by the letter of the statement, and vanishing in the limit where fluid
mechanics lives. The announcement therefore does not close the subject. It demonstrates, more clearly
than argument could, that the official alternatives can be discharged without bearing on turbulence
--- and that a formulation whose horns can be satisfied without touching its subject is the thing
that needs attention.

I would be glad to supply computational detail, or to discuss the formulation with anyone the Board
cares to involve.

\medskip
\noindent Yours faithfully,\\[2em]
\noindent Claes Johnson\\
\noindent KTH Royal Institute of Technology

\vspace{1.5em}
\begin{thebibliography}{9}
\bibitem{fujitakato} H.~Fujita and T.~Kato, \emph{On the Navier--Stokes initial value problem~I},
Arch.\ Rational Mech.\ Anal.\ \textbf{16} (1964) 269--315.

\bibitem{kato} T.~Kato, \emph{Strong $L^p$-solutions of the Navier--Stokes equation in
$\mathbb{R}^m$, with applications to weak solutions}, Math.\ Z.\ \textbf{187} (1984) 471--480.

\bibitem{kolmogorov41} A.~N.~Kolmogorov, \emph{The local structure of turbulence in incompressible
viscous fluid for very large Reynolds numbers}, Dokl.\ Akad.\ Nauk SSSR \textbf{30} (1941) 301--305.

\bibitem{obukhov41} A.~M.~Obukhov, \emph{On the distribution of energy in the spectrum of turbulent
flow}, Bull.\ Acad.\ Sci.\ USSR, Geog.\ Geophys.\ Ser.\ \textbf{5} (1941) 453--466.

\bibitem{kolmogorov62} A.~N.~Kolmogorov, \emph{A refinement of previous hypotheses concerning the
local structure of turbulence in a viscous incompressible fluid at high Reynolds number}, J.\ Fluid
Mech.\ \textbf{13} (1962) 82--85.

\bibitem{obukhov62} A.~M.~Obukhov, \emph{Some specific features of atmospheric turbulence}, J.\
Geophys.\ Res.\ \textbf{67} (1962) 3011--3014.

\bibitem{frischparisi} U.~Frisch and G.~Parisi, \emph{On the singularity structure of fully developed
turbulence}, in M.~Ghil, R.~Benzi and G.~Parisi (eds.), \emph{Turbulence and Predictability in
Geophysical Fluid Dynamics and Climate Dynamics}, North-Holland, 1985.

\bibitem{musci} B.~Musci, B.~Dubrulle, J.~LeBris, D.~Geneste, P.~Braganca, J.-M.~Foucaut, C.~Cuvier
and A.~Cheminet, \emph{Experimental measurement of the vorticity--strain alignment around extreme
energy transfer events}, J.\ Fluid Mech.\ \textbf{1016} (2025) A53.

\bibitem{iyer} K.~P.~Iyer, T.~D.~Drivas, G.~L.~Eyink and K.~R.~Sreenivasan, \emph{Turbulence without
walls: whither the zeroth law of turbulence?}, Phys.\ Rev.\ Lett.\ \textbf{135} (2025) 134001.

\bibitem{taylor} G.~I.~Taylor, \emph{Statistical theory of turbulence}, Proc.\ R.\ Soc.\ Lond.\ A
\textbf{151} (1935) 421--444.

\bibitem{eyink} G.~L.~Eyink, \emph{Energy dissipation without viscosity in ideal hydrodynamics~I.
Fourier analysis and local energy transfer}, Physica~D \textbf{78} (1994) 222--240.

\bibitem{dls} C.~De~Lellis and L.~Sz\'ekelyhidi~Jr., \emph{Dissipative Euler flows and Onsager's
conjecture}, J.\ Eur.\ Math.\ Soc.\ \textbf{16} (2014) 1467--1505.

\bibitem{bdlis} T.~Buckmaster, C.~De~Lellis, P.~Isett and L.~Sz\'ekelyhidi~Jr., \emph{Anomalous
dissipation for $1/5$-H\"older Euler flows}, Ann.\ of Math.\ \textbf{182} (2015) 127--172.

\bibitem{duchonrobert} J.~Duchon and R.~Robert, \emph{Inertial energy dissipation for weak solutions
of incompressible Euler and Navier--Stokes equations}, Nonlinearity \textbf{13} (2000) 249--255.

\bibitem{fehn} N.~Fehn, M.~Kronbichler and G.~Lube, \emph{From anomalous dissipation through Euler
singularities to stabilized finite element methods for turbulent flows}, Flow Turbul.\ Combust.\
\textbf{115} (2025) 347--388.

\bibitem{triple} J.~Kronborg and J.~Hoffman, \emph{The triple decomposition of the velocity gradient
tensor as a standardized real Schur form}, Phys.\ Fluids \textbf{35} (2023) 031703.

\bibitem{energystab} J.~Hoffman, \emph{Energy stability analysis of turbulent incompressible flow based
on the triple decomposition of the velocity gradient tensor}, Phys.\ Fluids \textbf{33} (2021) 081707.

\bibitem{hjv} J.~Hoffman, J.~Jansson and R.~Vilela de Abreu, \emph{Adaptive modeling of turbulent
flow with residual based turbulent kinetic energy dissipation}, Comput.\ Methods Appl.\ Mech.\
Engrg.\ \textbf{200} (2011) 2758--2767.

\bibitem{hjja} J.~Hoffman, J.~Jansson, N.~Jansson and R.~Vilela de Abreu, \emph{Towards a
parameter-free method for high Reynolds number turbulent flow simulation based on adaptive finite
element approximation}, Comput.\ Methods Appl.\ Mech.\ Engrg.\ \textbf{288} (2015) 60--74.

\bibitem{unicorn} J.~Hoffman, J.~Jansson \emph{et al.}, \emph{Unicorn: Parallel adaptive finite
  element simulation of turbulent flow and fluid--structure interaction}, Comput.\ Fluids, 2013.
\bibitem{massparallel} N.~Jansson, J.~Hoffman and J.~Jansson, \emph{Framework for massively parallel
  adaptive finite element computational fluid dynamics}, SIAM J.\ Sci.\ Comput., 2012.
\bibitem{fenicshpc} J.~Hoffman, J.~Jansson and N.~Jansson, \emph{FEniCS-HPC: Automated predictive
  high-performance finite element computing with aerodynamics applications}, Lecture Notes in
  Computer Science, Springer, 2016.
\bibitem{fenicsmulti} J.~Hoffman, J.~Jansson \emph{et al.}, \emph{FEniCS-HPC: Coupled multiphysics
  in computational fluid dynamics}, Lecture Notes in Computer Science, Springer, 2017.
\bibitem{paralleltime} J.~Jansson and J.~Hoffman, \emph{Direct FEM parallel-in-time computation of
  turbulent flow}, 2017.
\bibitem{dlrf11} J.~Jansson, J.~Hoffman, N.~Jansson and R.~Vilela de Abreu, \emph{Time-resolved
adaptive FEM simulation of the DLR-F11 aircraft model at high Reynolds number}, 2014.

\bibitem{newapproach} J.~Hoffman and C.~Johnson, \emph{A new approach to computational turbulence
modeling}, Comput.\ Methods Appl.\ Mech.\ Engrg.\ \textbf{195} (2006) 2865--2880.

\bibitem{openai} OpenAI, \emph{Finite time blowup for Navier--Stokes}, preprint, 2026.
\newblock \url{https://cdn.openai.com/pdf/32d9f210-8b73-45e0-91bc-82a30aef8a9a/navier-stokes.pdf}

\bibitem{buckmaster} T.~Buckmaster, \emph{Statement}, 2026.
\newblock \url{https://cims.nyu.edu/~tristanb/statement.pdf}

\bibitem{civ} P.~Constantin, M.~Ignatova and V.~Vicol, \emph{Regularity of asymptotically
axisymmetric solutions to the 3D Navier--Stokes equations with analytic forcing},
arXiv:2609.20803, 2026.

\bibitem{onsager} L.~Onsager, \emph{Statistical hydrodynamics}, Nuovo Cimento \textbf{6},
Suppl.~2 (1949) 279--287.

\bibitem{fefferman} C.~Fefferman, \emph{Existence and smoothness of the Navier--Stokes equation},
Official Problem Description, Clay Mathematics Institute, 2000.
\bibitem{leray} J.~Leray, \emph{Sur le mouvement d'un liquide visqueux emplissant l'espace},
Acta Math.\ \textbf{63} (1934) 193--248.
\bibitem{cet} P.~Constantin, W.~E and E.~S.~Titi, \emph{Onsager's conjecture on the energy
conservation for solutions of Euler's equation}, Comm.\ Math.\ Phys.\ \textbf{165} (1994) 207--209.
\bibitem{isett} P.~Isett, \emph{A proof of Onsager's conjecture}, Ann.\ of Math.\ \textbf{188}
(2018) 871--963.
\bibitem{cfd} J.~Hoffman and C.~Johnson, \emph{Computational Turbulent Incompressible Flow},
Applied Mathematics: Body and Soul, vol.~4, Springer, 2007.
\bibitem{dalembert} J.~Hoffman and C.~Johnson, \emph{Resolution of d'Alembert's paradox},
J.\ Math.\ Fluid Mech.\ \textbf{12} (2010) 321--334.
\bibitem{flight} J.~Hoffman, J.~Jansson and C.~Johnson, \emph{New theory of flight},
J.\ Math.\ Fluid Mech.\ \textbf{18} (2016) 219--241.
\bibitem{drag} J.~Hoffman, \emph{Efficient computation of mean drag for the subcritical flow past a
circular cylinder using general Galerkin G2}, Int.\ J.\ Numer.\ Meth.\ Fluids \textbf{59} (2009)
1241--1258.
\bibitem{lions} J.-L.~Lions, \emph{Quelques m\'ethodes de r\'esolution des probl\`emes aux
limites non lin\'eaires}, Dunod, Paris, 1969.

\bibitem{blog} C.~Johnson, posts under the label \emph{clay problem},
\url{https://claesjohnson.blogspot.com/search/label/clay\%20problem}.
\end{thebibliography}
\end{document}
